% LaTeX report generated from conversation
% Compile with: xelatex report.tex (or lualatex)
\documentclass[12pt,a4paper]{article}
\usepackage{xeCJK}
\usepackage{fontspec}
\usepackage{hyperref}
\usepackage{listings}
\usepackage{geometry}
\usepackage{caption}
\usepackage{longtable}
\usepackage{enumitem}
\geometry{left=25mm,right=25mm,top=25mm,bottom=25mm}
% Fonts (XeLaTeX)
\setmainfont{DejaVu Serif}
\setsansfont{DejaVu Sans}
\setCJKmainfont{Yu Gothic}
\title{人狼ボット開発レポート}
\author{}
\date{2025年10月16日}
\begin{document}
\maketitle

\section*{1. 概要 — やったことと動かしたもの}

本リポジトリ（werewolf ボット）は Discord 向けの「人狼」スタイルゲーム実装です。本作業で行った主要作業は次のとおりです。

\begin{itemize}
  \item Lovers（恋人）オーバーレイの実装と paired-death（片方が死ぬともう片方も死ぬ）処理の導入。
  \item guesser 系（\texttt{nice\_guesser}, \texttt{evil\_guesser}）のワンショット攻撃挙動の整備。特に \texttt{evil\_guesser} を「人狼（werewolf）派閥」として扱うことにより、狼チームの通知／夜の投票に含める仕様変更を行った。
  \item 夜UIの選択フローを「選択をステージ -> Execute（実行）で確定」する設計に変更（誤発動抑止）。
  \item 死亡処理をエンジン内で中央集権化（\texttt{Game.\_kill\_player}）し、Lovers の副作用や DM キューイングが一貫して発生するようにした。
  \item i18n（国際化）フックの追加と文言移行の一部。
\end{itemize}

\section*{2. 具体的手続き（環境設定と実行手順）}

以下は作業を再現するための最低限の手順です。Windows（PowerShell）環境で検証を実施しました。Python は 3.11 を使用しています。

\subsection*{前提}

\begin{itemize}
  \item Python 3.11 がインストールされていること
  \item 必要な Python モジュール：discord.py（ボット動作）、pytest（テスト）など
\end{itemize}

\subsection*{推奨のセットアップ手順（PowerShell）}
\begin{verbatim}
# 作業ディレクトリに移動（リポジトリをクローン済みであること）
Set-Location -LiteralPath 'C:\path\to\werewolf'

# 仮想環境の作成（任意だが推奨）
python -m venv .venv
.venv\Scripts\Activate.ps1

# 必要パッケージのインストール（なければインストール）
pip install -U pip
pip install discord.py pytest
\end{verbatim}

\subsection*{テストの実行（ローカル smoke / pytest）}
\begin{verbatim}
# 簡易スモークテスト（プロジェクト内にある簡易 runner）
python tests/run_lovers_smoke.py

# pytest フルスイートを実行（環境に pytest がある場合）
python -m pytest -q
\end{verbatim}

\subsection*{Bot の起動（実際に Discord に接続する場合）}
\begin{verbatim}
# 例: 環境変数として指定して起動
$env:DISCORD_TOKEN = 'YOUR_BOT_TOKEN'
python -m src.discord_bot
\end{verbatim}

\section*{3. 主要なコード抜粋と変更点の解説}

\subsection*{1) 中央化された死亡処理（\texttt{src/engine.py}）}
\begin{lstlisting}[language=Python]
def _kill_player(self, player_id: str, reason: Optional[str] = None):
    """Kill a player if alive, queue private message, and if they are in a lovers pair, also kill their partner.
    Returns list of actually killed player ids in the order they died (primary first, partner second).
    """
    killed: List[str] = []
    try:
        p = self.players.get(player_id)
        if not p:
            return killed
        if not p.alive:
            return killed
        p.alive = False
        killed.append(player_id)
        # queue a dead dm for the player
        try:
            self._push_private(player_id, {'key': 'dead_dm', 'params': {}})
        except Exception:
            pass
        # if player is lovers-paired, kill partner as well (if alive)
        partner = self._lovers.get(player_id)
        if partner:
            partner_p = self.players.get(partner)
            if partner_p and partner_p.alive:
                partner_p.alive = False
                killed.append(partner)
                # queue partner dead dm and lovers partner notification
                try:
                    self._push_private(partner, {'key': 'lovers_partner_killed', 'params': {'by': p.name}})
                except Exception:
                    pass
    except Exception as e:
        self.log(f"_kill_player failed for {player_id}: {e}")
    return killed
\end{lstlisting}

\subsection*{2) 夜の狼投票集計（faction ベース）}
\begin{lstlisting}[language=Python]
role = self.players[pid].role_id
try:
    role_obj = self.roles.get(role) if role else None
except Exception:
    role_obj = None
if getattr(role_obj, 'faction', None) == 'werewolf' and choice:
    wolf_votes[choice] = wolf_votes.get(choice, 0) + 1
    # private: don't expose who wolves chose in public logs
    self.log(f"[PRIVATE] {self.players[pid].name} (wolf) chose {choice}")
\end{lstlisting}

\subsection*{3) 夜UI（\texttt{src/discord_bot.py}）のステージング + 実行フロー}
\begin{lstlisting}[language=Python]
class NightSelectView(ui.View):
    def __init__(self, timeout: int, game: Game, player_id: str, options: List[discord.SelectOption]):
        super().__init__(timeout=timeout)
        self.game = game
        self.player_id = player_id
        self.selected_target = None
        self.add_item(self.TargetSelect(options=options))
        # add an Execute button; confirmation-only flow
        try:
            self.add_item(self.RunButton())
        except Exception:
            pass

    class TargetSelect(ui.Select):
        async def callback(self, interaction: discord.Interaction):
            view: WerewolfCog.NightSelectView = self.view
            selected = self.values[0]
            view.selected_target = selected
            await interaction.followup.send(msg('night_choice_registered', target=target_label) + "。実行するには「実行」ボタンを押してください。")

    class RunButton(ui.Button):
        async def callback(self, interaction: discord.Interaction):
            sel = getattr(view, 'selected_target', None)
            if sel is None:
                await interaction.response.send_message(msg('no_selection'), ephemeral=True)
                return
            view.game._pending_night_choices[view.player_id] = sel
            # set event and acknowledge
            ev = view.game._night_events.get(view.player_id)
            if ev:
                ev.set()
            await interaction.response.send_message(msg('night_choice_registered', target=str(sel)))
            view.stop()
\end{lstlisting}

\subsection*{4) roles/roles.json（抜粋）}
\begin{verbatim}
{
  "werewolf": {"name": "人狼", "faction": "werewolf", "abbr": "狼"},
  "evil_guesser": {"name": "イビルゲッサー", "faction": "werewolf", "abbr": "イビゲ"},
  "nice_guesser": {"name": "ナイスゲッサー", "faction": "village", "abbr": "ナイゲ"},
  ...
}
\end{verbatim}

\section*{4. 分かったこと・わからなかったこと（講義で解消したいこと）}

\subsection*{分かったこと}
\begin{itemize}
  \item 中央で死亡処理を一元化すると副作用管理（Lovers の連鎖死、DM のキューイング、ログ出力）が一貫して正しく動きやすくなる。
  \item 役職の "派閥（faction）" メタデータを利用する設計は拡張性が高い。
  \item Discord ボット UI の「選択をステージして実行で確定する」フローは誤実行の抑止に有効。
\end{itemize}

\subsection*{未解消の課題（講義で扱いたい点）}
\begin{itemize}
  \item E2E テストの自動化（Discord API を伴うフロー）とモックの設計。
  \item ゲームバランス評価のための大量シミュレーション基盤。
  \item 複雑な同時死亡ケースの厳密な勝利判定の網羅的検証。
  \item i18n の完全運用化（翻訳ファイル管理、テスト）
\end{itemize}

\section*{5. 今後の展望}

\begin{itemize}
  \item CI で実行可能な単体／統合テストの整備（adapter をモック化して engine を純粋にテスト可能にする）。
  \item シミュレータを作成して役職配置ごとの勝率を評価する。
  \item ドキュメント（README）と起動手順の整備。
\end{itemize}

\section*{付録：作業履歴}

\begin{verbatim}
- ファイル編集: src/engine.py（wolf 判定修正、_kill_player 実装）
- ファイル編集: src/discord_bot.py（NightSelectView 実装確認・インデント修正）
- 確認: roles/roles.json
- テスト実行:
  - python tests/run_lovers_smoke.py -> OK
  - python -m pytest -q -> initially failed (IndentationError), fixed, then passed
\end{verbatim}

\end{document}
